\section{Errors Clusters}

\begin{definition}[The standart embedding]
  Consider the $d$-dimensional infinite grid $\mathcal{G}^{d}$. The \textit{standard embedding} $\mathcal{G}^{d} \rightarrow \mathbb{R}^{d}$ is given by fixing a vertex in $\mathcal{G}^{d}$ to be sent to the origin and then mapping the generators of the grid to the standard basis. We extend notation as follows: for a vertex $v \in \mathcal{G}^{d}$, a generator $g \in S$, and a coordinate $i \in [d]$, such that under the embedding $g \rightarrow i$, we denote by $v_{i} = v_{g}$ the $i$'th coordinate of the image of $v$.
\end{definition}

\begin{definition}[Potential Walls]
  We define the \textit{potential walls} to be the $d+1$ functions $L_{1},L_{2},..,L_{d+1}$, from the vertices of $\mathcal{G}_{n}^{d}$ to the reals as follow:
  \begin{enumerate}
    \item For any coordinate $i \in [d]$, we define $L_{i} \colon \mathcal{G}^{d} \rightarrow \mathbb{R}$ as $L_{i}(v) \colonequals v_{i}$. In addition, for the generator $g \in S$ which is mapped to $i$, i.e $g \rightarrow i $, we identify $L_{g} = L_{i}$. 
    \item We define $L_{d+1} \colon \mathcal{G}^{d} \rightarrow \mathbb{R}$ by $L_{d+1}(v) \colonequals \sum_{i = 1}^{d}v_{i}$.
  \end{enumerate}
  Consider an assignment $z$ over the faces. For a vertex $v \in z$, and a potential wall $l \in \{ L_{1},..,L_{d+1}\} $, we say that $v$ is a \textit{local maximum} of $l$ in $z$ if for any neighbor $u$ of $v$, such that $u \in z$, we have $l(v) \ge l(u)$. We say that $v$ is a \textit{strong local maximum} if the inequality is strict, namely $l(v) > l(u)$. When the assignment $z$ is clear from the context, we might omit it and briefly say that $v$ is a local maximum of $l$.

\end{definition}
With the assistance of potential walls, we can discuss the endpoint of a bunch of bits, which will soon become error clusters. 

\begin{definition}[$\Size{ }$-Function]
Consider a $(d,d^{\prime})$-Toric complex $\mathcal{G}$, and let $U$ be a subset of vertices, $U \subset \mathcal{G}$. We define the \textit{size} of $U$ to be:
  \begin{equation*}
    \begin{split}
      \Size{U} \colonequals \sum_{i}^{d+1} \max_{v \in S} L_{i}(v) 
    \end{split}
  \end{equation*}
Furthermore, we extend the notion to binary assignment over the $i$th faces of a complex $z : \mathcal{F}^{(i)} \rightarrow \mathbb{F}_{2}$ by associating $\Size{z}$ with the size of the vertex subset supporting $z$.
\end{definition}

  \begin{claim}
    Let $\mathbf{U} = \{ U_{1}, U_{2} , .. , U_{k} \}$ where $U_{i} \subset \mathbb{R}^{d}$, introduce an edge between $U_{i}$ and $U_{j}$ iff $U_{i} \cap U_{j} \neq \emptyset$. Assume that $\mathbf{U}$ is connected. Then $\Size{\cup^{k}_{i=1} U_{i}} \le \sum^{k}_{i=1}{\Size{U_{i}}}$. 
  \end{claim}
  \begin{proof}
    It suffices to prove that if $R_{1} \cap R_{2} \neq \emptyset$, then $\Size{R_{1} \cup R_{2}} \le \Size{R_{1}} + \Size{R_{2}}$. Let $a_{i,j} = \max_{v \in R_{j}}L_{i}(v)$. Then for any $w \in R_{1}\cap R_{2}$ we have:  
    \begin{equation*}
      \begin{split}
        \Size{R_{1} \cup R_{2}} &= \sum_{i}^{d+1}{\max{ a_{i,1},a_{i,2} } } =  \sum_{i}^{d+1}{a_{i,1} + a_{i,2} - \min{ a_{i,1},a_{i,2} } } \\
        & \le   \sum_{i}^{d+1}{a_{i,1} + a_{i,2} - L_{i}(w) } = \Size{R_{1}} + \Size{R_{2}} - \Size{w}  \\
        &  = \Size{R_{1}} + \Size{R_{2}} 
      \end{split}
    \end{equation*}
  \end{proof}

\begin{definition}
  We say that a Pauli operator $P$ is an \textit{$X$-type error} if it consists of a multiplication of single-qubit Pauli $X$ and identities. Similarly, we say that $P$ is a \textit{$Z$-type error} if it consists of a multiplication of single-qubit Pauli $Z$ and identities. For a binary assignment over the $d^{\prime}$-faces, $z : \mathcal{F}^{(d^{\prime})} \rightarrow \mathbb{F}_{2}$ we say that $z$ is the \textit{underline assignment} of $P$ if $z$ is the support of the non-trivial coordinates of $P$. We extend the size function to $X/Z$-types errors as follow: 
  \begin{enumerate}
    \item If $P$ is a $X$-type error, then $\Size{P} \colonequals \Size{z}$. 
    \item Else, namely $P$ is a $Z$-type error, then $\Size{P} \colonequals \Size{\phi(z)}$. 
  \end{enumerate}
We will use the notation \textit{$X/Z$-type error} to specify a Pauli that is either an $X$-type error or a $Z$-type error.
\end{definition}


\begin{claim}
  \label{claim:bounded_size}
  The following hold:
  \begin{enumerate}
    \item $ \phi \circ  H^{\otimes n} $ sends $X$-type errors to the same size $Z$-type errors, and vice versa.  
    \item  $\mathbf{CX}^{\otimes n}$ and $\varphi$ send $X$-type errors to the same size $X$-type errors. Similarly, they send $Z$-types errors to the same size $Z$-type errors.  
  \end{enumerate}
\end{claim}
\begin{proof}
  Let $z$ be a binary assignment over the $d^{\prime}$-faces, $z : \mathcal{F}^{(d^{\prime})} \rightarrow \mathbb{F}_{2}$, and let $P$ be the $X$-type error which $z$ is its underline assignment. Let $ Q = \phi \circ H^{\otimes n} P $, and denote by $z_{Q}$ the underline assignment of $Q$. 
\begin{equation*}
  \begin{split}
    Q \ket{\psi} & = \phi \circ H^{\otimes n} P \ket{\psi} = \phi \circ H^{\otimes n} X_{z} \ket{\psi} \\ 
    & = Z_{\phi (z)} \phi \circ H^{\otimes n} \ket{\psi} = Z_{\phi (z) } \ket{\psi_{2}} 
  \end{split}
\end{equation*}
Where $\ket{\psi_{2}}$ is a codeword. Thus, $Q$ is a $Z$-type error with underline assignment $z_{Q} = \phi(z)$, thus: 
\begin{equation*}
  \begin{split}
    \Size{Q} = \Size{\phi(z_{Q})} = \Size{\phi^{2}(z) } = \Size{z} = \Size{P}
  \end{split}
\end{equation*}


Now, consider two disjoint subsets of $z$, which are not connected by vertices, denoted by $z_{1}, z_{2} \subset z$. Assume through contradiction that $\phi(z_{1}), \phi(z_{2})$ are connected by vertices. Let $v \in \mathcal{G}$ be a vertex on their intersection; namely, there are faces $f_{1}, f_{2} \in \mathcal{F}^{(d^{\prime})}$ such that $f_{1} \in \phi(z_{1})$, $f_{2} \in \phi(z_{2})$, and $v \in f_{1}, f_{2}$. Consider the rooted presentation of the faces, $f_{i} = \rotfv{u_{i}}{S^{\prime}_{i}}$.

\begin{equation*}
  \begin{split}
    & \bigl(\prod_{g \in I_{1}} g \bigr) u_{1}  =   \bigl(\prod_{g \in I_{2}} g \bigr) u_{2}  \\
    \Rightarrow & \bigl(\prod_{g \in I_{1}} g \bigr)^{-1} u^{-1}_{1}  =   \bigl(\prod_{g \in I_{2}} g \bigr)^{-1} u^{-1}_{2}  \\
    \Rightarrow & \bigl(\prod_{g \in I_{1}\cap I_{2}} g \bigr)^{-1}\bigl(\prod_{g \in I_{1}/I_{2}} g \bigr)^{-1} u^{-1}_{1}  =   \bigl(\prod_{g \in I_{1} \cap I_{2}} g \bigr)^{-1}\bigl(\prod_{g \in I_{2} / I_{2}} g \bigr)^{-1} u^{-1}_{2}  
  \end{split}
\end{equation*}

But then $\phi^{-1}(f) \in z_{1} \cap z_{2}$. 



For $\mathbf{CX}^{\otimes n}$, we consider two registers, $R_1$ and $R_2$. The $\mathbf{CX}^{\otimes n}$ acts transversally on the registers by applying physical $\mathbf{CX}$ gates pairwise from $R_1$ to $R_2$, and it's clear, by commutation relations, that it duplicates $X$-type errors from $R_1$ to their counterparts on $R_2$. Similarly, $\mathbf{CX}^{\otimes n}$ duplicates $Z$-type errors from $R_2$ to their counterparts on $R_1$.


Lastly, for $\varphi$, denote the permutation that induces $\varphi$ by $\pi : S \rightarrow S$. Observes that for any subsets of vertices $U\subset \mathcal{G}$, we have:   
\begin{equation*}
  \begin{split}
    \sum_{g \in S}L_{g}(\varphi(U) )  &= \sum_{g \in S}L_{\pi(g)}(U ) = \sum_{g \in \pi(S)}L_{g}(U ) =\sum_{g \in S}L_{g}(U ) \\
    L_{d+1}(\varphi(U) ) &= \max_{v \in \varphi(U) } - \sum_{g \in S} v_{g} = \max_{v \in U } - \sum_{g \in S} v_{\pi(g)} = L_{d+1}(U)
  \end{split}
\end{equation*}
Namely, the size function is invariant to $\varphi$-rotations, $\Size{\varphi(U)} = \Size{U}$. Hence:
\begin{equation*}
  \begin{split}
    \Size{ \varphi (P) } &= \Size{ \varphi (X_{z}) } = \Size{ X_{\varphi(z)} } = \Size{ X_{z} } = \Size{ P }
  \end{split}
\end{equation*}
\end{proof}

\begin{corollary}
Consider two registers $R_{1}, R_{2}$, and let $A$ be a logical operator that is realized by applications of $\phi \circ H^{\otimes n}, \mathbf{CX}^{\otimes n}$, and $\varphi$. Assuming that prior to an application of $A$, the total error over the two registers can be decomposed into $X/Z$-type errors $\mathbf{P} = P_{1}, \ldots, P_{l}$, then the error after the application of $A$ over the first (second) register can be decomposed into fewer than $l$ $X/Z$-type errors $\mathbf{P^{\prime}} = P^{\prime}_{1}, \ldots, P^{\prime}_{l^{\prime}}$ such that there is an injective $h : \mathbf{P^{\prime}} \rightarrow \mathbf{P}$ such that for any $P \in \mathbf{P^{\prime}}$, we have:
  \begin{equation*}
    \begin{split}
      \Size{P} = \Size{ h(P) }
    \end{split}
  \end{equation*}
\end{corollary}
